
\documentclass[11pt]{article}

\usepackage{amsmath,amssymb,amsthm,mathtools}
\usepackage[T1]{fontenc}
\usepackage{lmodern}
\usepackage[margin=1in]{geometry}
\usepackage{enumitem}
\usepackage{booktabs}
\usepackage{array}
\usepackage{longtable}
\usepackage{microtype}
\usepackage{xcolor}
\usepackage{hyperref}
\usepackage{xurl}

\hypersetup{colorlinks=true,linkcolor=blue,citecolor=blue,urlcolor=blue}
\urlstyle{same}

\newtheorem{theorem}{Theorem}[section]
\newtheorem{frontier}[theorem]{Frontier theorem}
\newtheorem{background}[theorem]{Background input}
\newtheorem{proposition}[theorem]{Proposition}
\newtheorem{lemma}[theorem]{Lemma}
\newtheorem{corollary}[theorem]{Corollary}
\newtheorem{remark}[theorem]{Remark}
\newtheorem{definition}[theorem]{Definition}
\newtheorem{heuristic}[theorem]{Heuristic}
\newtheorem{question}[theorem]{Open problem}

\newcommand{\sig}[2]{\sigma_{#1#2}}
\newcommand{\Lw}[1]{L_{#1}}
\newcommand{\Rw}[1]{R_{#1}}
\newcommand{\CanL}{[2,2,1]}
\newcommand{\CanR}{[1,4,4]}
\newcommand{\CanSeed}{\CanL/\CanR}
\newcommand{\CL}{\mathcal C_L}
\newcommand{\CR}{\mathcal C_R}
\newcommand{\Omix}{\mathcal O_{\mathrm{mix}}}
\newcommand{\Obr}{\mathcal O_{\mathrm{br}}}
\newcommand{\RhoR}{\rho_{R1}}
\newcommand{\SelCorr}{\mathrm{Sel}_{\mathrm{corr}}}
\newcommand{\SelStar}{\mathrm{Sel}^{\star}}

\title{Odd-$m$, $d=5$ as a one-cornered odometer:\\
current self-contained working manuscript}
\author{}
\date{March 2026}

\begin{document}
\maketitle

\begin{abstract}
This text is a self-contained working manuscript for the current odd-$m$,
$d=5$ proof strategy.
The proof is organized around the spine
\[
\begin{aligned}
&\text{front-end one-corner reduction}
\to
\text{exact odometer spine}\\
&\to
\text{selector surgery}
\to
\text{final stitching and globalization}.
\end{aligned}
\]
The front end is no longer described by historical package names.
Instead it is reduced to a finite chart and four theorem objects:
three shape-side row theorems
\[
T1,\ T2,\ T3
\]
and one canonical entry theorem
\[
T4.
\]
As of the present cleanup, the actual-needed tiny-repair core and the three
shape-side row theorems $T1,T2,T3$ and the canonical entry theorem $T4$ are
already closed.  The odometer block after entry is therefore a fixed
front-end input.
On the graph side, one full Hamilton branch is already modeled, and the
remaining selector-compatibility step has now been closed by an explicit
two-defect-set splice into the valid baseline package.  The remaining-color
transport theorem \(G2\) is already closed separately by cyclic conjugacy of the
closed branch.

The manuscript distinguishes carefully between three kinds of statements:
finite certificates, odd-uniform symbolic theorems, and explanatory heuristics.
At the present stage there is no live frontier theorem left inside the present
\(T0\)--\(T4,G1,G2\) slice.  The blocks still treated as imported rather
than reproved in this bundle are the post-entry odometer spine and the
stable graph-side/globalization inputs listed later in
Theorem~\ref{thm:odometer-spine} and Proposition~\ref{prop:stable-graph-side}.
\end{abstract}

\tableofcontents

\section{The global proof spine}

The intended theorem-level structure of the odd-$m$, $d=5$ paper is
\[
\text{front-end one-corner reduction}
\to
\text{exact odometer spine}
\to
\text{selector surgery}
\to
\text{final stitching}.
\]
At a slightly finer level the manuscript is organized as follows.
\begin{enumerate}[label=\textup{(\Roman*)},leftmargin=3.0em]
\item \textbf{Front-end one-corner reduction.}
      Reduce every unresolved front-end instance to the canonical bridge seed
      \(\CanSeed\) by a finite chart plus three shape-side row theorems.
\item \textbf{Canonical entry and exact odometer spine.}
      Show that the canonical seed has a unique live unresolved continuation,
      namely the alternate-$2$ entry
      \[
      E=(m-1,1,u_{\mathrm{source}},2),
      \]
      and then identify the resulting dynamics with the exact one-corner
      odometer corridor.
\item \textbf{Selector surgery.}
      Insert the already closed color-$4$ Hamilton branch into a colorwise
      selector package by a selector-compatibility theorem.  The remaining
      colors then follow by cyclic transport.
\item \textbf{Final stitching and globalization.}
      Use the colorwise graph package together with the explicit odometer block
      to produce the final odd-$m$, $d=5$ Hamilton decomposition.
\end{enumerate}

\begin{theorem}[actual-needed tiny-repair core]
\label{thm:T0-needed}
For every odd modulus $m>2$, the front-end uses of tiny repair that actually
occur below are already closed internally.
More precisely:
\begin{enumerate}[label=\textup{(\arabic*)},leftmargin=2.8em]
\item at
      \[
      \sig21=[4,1,4]/[1,2,2],
      \]
      every one-label one-gate or one-bit repair branch supported on
      \(L2,L3,R2,\) or \(R3\), with cocycle-defect choice in
      \(\{\mathrm{none},\mathrm{left},\mathrm{right},\mathrm{both}\}\),
      is non-Latin;
\item at the canonical bridge seed \(\CanSeed\), every one-label one-gate or
      one-bit repair branch supported on \(L3,R2,\) or \(R3\), with the same
      four cocycle-defect choices, is non-Latin.
\end{enumerate}
\end{theorem}

\begin{remark}[scope convention for ``after tiny repair'']
Throughout the front-end section, the phrase ``after tiny repair'' refers only
to the concrete removals supplied by Theorem~\ref{thm:T0-needed} at the relevant
seed.  It does not silently quantify over richer multi-class repair families.
A stronger standalone one-label repair theorem for the full artifact seed list
may still be desirable later, but it is not used below.
\end{remark}

\begin{theorem}[packaged full odd-$m$, $d=5$ theorem]
\label{thm:conditional-full}
Assume the imported post-entry odometer block from
Theorem~\ref{thm:odometer-spine} and the stable graph-side inputs from
Proposition~\ref{prop:stable-graph-side}.
Then the odd-$m$, $d=5$ system admits the desired global Hamilton structure.
\end{theorem}

The purpose of the present draft is not to pretend that every theorem in
Theorem~\ref{thm:conditional-full} is already proved.
Rather, the goal is to make the proof architecture mathematically honest:
closed blocks are stated as closed blocks, and the unresolved proof objects are
spelled out explicitly.

\section{Three languages and a methodological rule}

\begin{remark}[support language, grid language, odometer language]
Three languages appear repeatedly and should not be confused.
\begin{enumerate}[label=\textup{(\arabic*)},leftmargin=2.8em]
\item \textbf{Support language.}
      This records only which local families are present.
      It is a pruning language.
\item \textbf{Grid language.}
      This records actual endpoint words, equivalently the grid cell
      \(\sig{i}{j}\).
      It is the main language of shape-side row theorems.
\item \textbf{Odometer language.}
      This records the explicit post-entry coordinates
      \((q,w,u,\Theta)\).
      It is the main language after the canonical entry theorem.
\end{enumerate}
The key rule is:
\emph{support is a pruning invariant, not a classification invariant.}
So support can help isolate a theorem target, but it cannot by itself separate
all rows that the grid language can separate.
\end{remark}

\section{Finite front-end chart}

\subsection{\texorpdfstring{The \(3\times 3\)}{The 3x3} bridge-position grid}

Define
\[
L_1=[1,4,4],\qquad L_2=[4,1,4],\qquad L_3=[4,4,1],
\]
and
\[
R_1=[1,2,2],\qquad R_2=[2,1,2],\qquad R_3=[2,2,1].
\]
The noncanonical rows are arranged in the grid
\[
\sig{i}{j}=L_i/R_j,\qquad 1\le i,j\le 3.
\]
The interpretation is:
\begin{itemize}[leftmargin=2.4em]
\item \(i\) records the position of the bridge symbol on the left endpoint;
\item \(j\) records the position of the bridge symbol on the right endpoint.
\end{itemize}

\begin{proposition}[finite seed certificate]
\label{prop:finite-chart}
The filtered front-end seed set contains exactly eleven raw rows.
Under the accepted symmetries, these reduce to five classes
\[
C0,\ C1,\ C2,\ C3,\ C4.
\]
Exactly one class is an oriented bridge class, namely \(C0\), and it contains
the canonical bridge seed
\[
\CanSeed=\CanL/\CanR.
\]
The centered-bridge classes are exactly
\[
C1\cup C3\cup C4,
\]
and the residual nonbridge class outside that layer is exactly
\[
C2=\sig33=[4,4,1]/[2,2,1].
\]
\end{proposition}

\begin{proposition}[bridge normalization]
\label{prop:bridge-normalization}
Every oriented bridge row is equivalent under the accepted symmetries to the
canonical seed
\[
\CanSeed=[2,2,1]/[1,4,4].
\]
\end{proposition}

\begin{definition}[changed-family support]
For each surviving row \(\sigma\), let
\[
\Sigma(\sigma)\subseteq\{L1,L2,L3,R1,R2,R3\}
\]
denote the changed-family support.
On the noncanonical \(3\times 3\) grid the support factorization is
\[
\Sigma(\sig{i}{j})
=
\{L2,L3,R1,R3\}
\cup (\{L1\}\ \text{if } i=1)
\cup (\{R2\}\ \text{if } j\le 2).
\]
\end{definition}

\begin{remark}[the exceptional row]
The row
\[
\sig11=[1,4,4]/[1,2,2]
\]
is exceptional for a finite reason:
it is the unique bridge-first row in which the left bridge is still leading.
In the bridge-position grid it is the double-leading corner.
\end{remark}

\subsection{Prototype reduction on the shape side}

\begin{proposition}[two-prototype reduction]
\label{prop:two-prototypes}
Under the accepted symmetries, the centered-bridge layer decomposes into two
orbits:
\[
\Omix=
\{\sig21,\ \sig23,\ \sig12,\ \sig32\},
\qquad
\{\sig22\}.
\]
Thus the whole centered-bridge burden reduces to one mixed prototype
\(\sig21=[4,1,4]/[1,2,2]\) and one double-centered core \(\sig22\).
\end{proposition}

\begin{corollary}[shape-side compression]
\label{cor:shape-side-compression}
Modulo Theorem~\ref{thm:T0-needed} and bridge normalization, the front-end shape side
reduces to exactly three row theorem objects:
\[
T1:\ \sig21\rightsquigarrow \sig31,\qquad
T2:\ \sig22\text{ closes},\qquad
T3:\ \sig33\text{ closes}.
\]
\end{corollary}

\section{Closed front-end row theorems}
\label{sec:frontier}

\begin{theorem}[T1: first-column companion theorem]
\label{fr:T1}
After the \(\sig21\)-repair branches from Theorem~\ref{thm:T0-needed} are
removed, every live unresolved immediate continuation of
\[
\sig21=[4,1,4]/[1,2,2]
\]
is represented by the bridge companion
\[
\sig31=[4,4,1]/[1,2,2].
\]
\end{theorem}

\begin{theorem}[T2: double-centered closure]
\label{fr:T2}\label{thm:T2-closed}
For every odd modulus \(m>2\), no row obtained from
\[
\sig22=[4,1,4]/[2,1,2]
\]
by a one-gate or one-bit one-label repair, with cocycle-defect choice in
\(\{\mathrm{none},\mathrm{left},\mathrm{right},\mathrm{both}\}\), is Latin.
Equivalently, after the artifact-surfaced tiny-repair scope is exhausted, the
center cell \(\sig22\) does not remain as a live unresolved front-end row.
\end{theorem}

\begin{theorem}[T3: repeated-end corner closure]
\label{fr:T3}\label{thm:T3-closed}
For every odd modulus \(m>2\), no row obtained from
\[
\sig33=[4,4,1]/[2,2,1]
\]
by a one-gate or one-bit one-label repair, with cocycle-defect choice in
\(\{\mathrm{none},\mathrm{left},\mathrm{right},\mathrm{both}\}\), is Latin.
Equivalently, after the artifact-surfaced tiny-repair scope is exhausted, the
repeated-end corner \(\sig33\) does not remain as a live unresolved front-end row.
\end{theorem}

\begin{remark}[proof packets for \(T2\) and \(T3\)]
The odd-\(m\ge 5\) obstruction proofs are extracted in the companion notes
\path{d5_228_T2_centered_core_obstruction_note.tex} and
\path{d5_225_T3_four_collision_families_note.tex}.
The remaining one-label \(m=3\) cases are exhausted by the finite scans recorded
in \path{d5_229_T2_T3_m3_one_label_scan.py} and its saved summaries.
\end{remark}

\begin{proposition}[canonical support certificate]
\label{prop:canonical-support}
At the canonical bridge seed
\[
\CanSeed=\CanL/\CanR,
\]
the support is exactly
\[
\Sigma(\CanSeed)=\{L3,R1,R2,R3\}.
\]
\end{proposition}

\begin{theorem}[T4: canonical entry theorem]
\label{fr:T4}
At the canonical bridge seed \(\CanSeed\), after the three dead repair source
families from Theorem~\ref{thm:T0-needed}---namely \(L3\), \(R2\), and
\(R3\)---are removed, the unique live immediate source family is \(R1\).
Its canonical two-step normalization in the shifted section coordinates
\[
(Q,W,U,\Theta)=(q-1,\ w+1,\ u-1,\ 2)
\]
is the alternate-$2$ entry
\[
E=(m-1,1,u_{\mathrm{source}},2).
\]
\end{theorem}

\begin{theorem}[front-end one-corner reduction]
\label{thm:conditional-front-end}
Every unresolved odd-$m$, $d=5$ front-end instance reduces to the exact
one-corner channel with canonical seed \CanSeed and entry state
\[
E=(m-1,1,u_{\mathrm{source}},2).
\]
\end{theorem}

\begin{proof}
By Proposition~\ref{prop:finite-chart}, every unresolved front-end instance is
represented by one of the finitely many rows in the chart.
Proposition~\ref{prop:bridge-normalization} moves every oriented bridge row to
the canonical seed \CanSeed.  Corollary~\ref{cor:shape-side-compression}
reduces the noncanonical shape side to the theorem objects \(T1,T2,T3\),
and Theorems~\ref{fr:T1}--\ref{fr:T4} close those row objects and identify
the unique live bridge continuation with the entry state \(E\).  Hence every
unresolved front-end instance reduces to the exact one-corner channel with
canonical seed \CanSeed and entry state \(E\).
\end{proof}

\section{\texorpdfstring{The closure package for \(T1\)}{The closure package for T1}}

What used to be the live front-end proof search is now a closed package.
We keep the strip analysis here because it explains why Theorem~\ref{fr:T1} is
short.

\subsection{Unique surviving branch}

\begin{proposition}[source support for \(T1\)]
\label{prop:T1-source-support}
At the mixed prototype
\[
\sig21=[4,1,4]/[1,2,2]
\]
one has
\[
\Sigma(\sig21)=\{L2,L3,R1,R2,R3\}.
\]
By Theorem~\ref{thm:T0-needed}, every immediate branch except the
\(R1\)-branch is removed.
\end{proposition}

So \(T1\) is not a many-branch problem.
It is a unique surviving-branch problem.

\subsection{Two strips and two local raw targets}

Define the centered-left strip and centered-right strip by
\[
\CL=\{\sig21,\sig22,\sig23\},
\qquad
\CR=\{\sig12,\sig22,\sig32\}.
\]
With Theorems~\ref{thm:T2-closed} and \ref{thm:T3-closed} recorded as
closed obstruction cells, \(T1\) reduces to excluding the dangerous targets
on these two strips.

\begin{proposition}[left-strip case split]
\label{prop:left-strip-case-split}
Relative to the source right word \([1,2,2]\), the centered-left strip has
exactly three cases:
\[
\sig21 \leftrightarrow \Delta_R=\varnothing,\qquad
\sig22 \leftrightarrow \Delta_R=\{1,2\},\qquad
\sig23 \leftrightarrow \Delta_R=\{1,3\}.
\]
In particular,
\[
\sig23=[4,1,4]/[2,2,1]
\]
is the unique centered-left case with disconnected change-set \(\{1,3\}\), the
unique centered-left case with terminal right-slot change, and the unique
centered-left case with \(R2\) absent.
\end{proposition}

\begin{proposition}[current \(T1\) reduction after strip closure]
\label{prop:T1-current-reduction}
Assume Theorems~\ref{thm:T2-closed} and \ref{thm:T3-closed}.
Then \(T1\) has no remaining live raw local target.
The left-strip target that the surviving immediate \(R1\)-image out of
\(\sig21\) does not realize \(\sig23\) is closed by
Theorem~\ref{thm:sigma23-closed}, and the normalized centered-right target that
it does not realize \(\sig32\) is closed by
Theorem~\ref{thm:sigma32-closed} below.
\end{proposition}

\subsection{The \texorpdfstring{$\sig23$}{sigma23} side is now closed}

The bookkeeping reductions from the earlier notes remain useful, but the live
frontier theorem on this side is no longer open.  Once the surfaced bundle
builder is admitted, the surviving immediate \(R1\)-family out of \(\sig21\)
can be read directly and its second-step comparison is explicit.

\begin{proposition}[equivalent forms of \(\sig23\) exclusion]
\label{prop:sigma23-equivalences}
For the surviving immediate \(R1\)-image \(\RhoR\) out of \(\sig21\), each of the
following suffices to exclude \(\sig23\).
\begin{enumerate}[label=\textup{(\arabic*)},leftmargin=2.8em]
\item \(\Delta_R(\RhoR)\subseteq\{1,2\}\);
\item the terminal right slot is unchanged;
\item the changed right slots form a connected set;
\item the immediate update is an elementary rewrite rooted at the acted edge
      \(\{1,2\}\);
\item the immediate update is supported on the closed star of the acted edge.
\end{enumerate}
\end{proposition}

\begin{theorem}[surfaced-bundle closure of the \(\sig23\) side]
\label{thm:sigma23-closed}
For every modulus \(m>2\), the actual second step of the surviving immediate
\(R1\)-family out of
\[
\sig21=[4,1,4]/[1,2,2]
\]
agrees with the \(\sig22=[4,1,4]/[2,1,2]\) prototype and never with the
\(\sig23=[4,1,4]/[2,2,1]\) prototype.
Equivalently, for the tiny-repair-normalized surviving immediate \(R1\)-image
\(\RhoR\) one has
\[
\Delta_R(\RhoR)=\{1,2\}.
\]
In particular the left-strip burden inside \(T1\) is closed.
\end{theorem}

\begin{proof}
Fix a modulus \(m>2\), color \(0\), and the standard normalization
\(w_0=s_0=0\).
Write \(\partial_i\) for the coordinate step that adds \(+1\) in direction
\(i\) modulo \(m\).
In the surfaced builder, the common right \(R1\)-source family for
\(\sig21,\sig22,\sig23\) is exactly
\[
S_R=\{x:\ \mathrm{layer}(x)=1,\ q(x)=m-1,\ w(x)=0,\ s(x)\neq 0\}.
\]
Thus all three rows are compared on the same source family.

For \(\sig21\), the right word is \([1,2,2]\), so the first two right-side
steps on \(S_R\) are
\[
x\mapsto \partial_1(x)\mapsto \partial_2\partial_1(x).
\]
For \(\sig22\), the right word is \([2,1,2]\), so the first two right-side
steps are
\[
x\mapsto \partial_2(x)\mapsto \partial_1\partial_2(x).
\]
For \(\sig23\), the right word is \([2,2,1]\), so the first two right-side
steps are
\[
x\mapsto \partial_2(x)\mapsto \partial_2^2(x).
\]

The surfaced common module defines each \(\partial_i\) as coordinatewise
\(+1\) in a single axis modulo \(m\).  Therefore the coordinate steps
commute:
\[
\partial_2\partial_1=\partial_1\partial_2.
\]
Thus the actual second step of \(\sig21\) on the surviving \(R1\)-family agrees
exactly with the \(\sig22\) prototype.

On the other hand, every \(x\in S_R\) satisfies \((q(x),w(x))=(m-1,0)\).
Therefore
\[
\partial_2\partial_1(x):(q,w)=(0,1),\qquad
\partial_2^2(x):(q,w)=(m-1,2).
\]
These are distinct for every \(m>2\).  Hence the actual second step of
\(\sig21\) does not realize \(\sig23\).
The description of \(\Delta_R\) now follows from
Proposition~\ref{prop:left-strip-case-split}.
\end{proof}

\subsection{The \texorpdfstring{$\sig32$}{sigma32} side is also closed}

The remaining centered-right target is the repeated-end option
\[
\sig32=[4,4,1]/[2,1,2].
\]
Here the common comparison family is on the left rather than on the right.
The first two left-side steps are again explicit in the surfaced builder.

\begin{theorem}[surfaced-bundle closure of the \(\sig32\) side]
\label{thm:sigma32-closed}
For every modulus \(m>2\), the normalized surviving immediate \(R1\)-image out
of
\[
\sig21=[4,1,4]/[1,2,2]
\]
does not realize the dangerous centered-right target
\[
\sig32=[4,4,1]/[2,1,2].
\]
More precisely, on the common left family
\[
S_L=\{x:\ \mathrm{layer}(x)=1,\ q(x)=m-1,\ w(x)=m-1,\ s(x)\neq 0\},
\]
the actual two-step left continuation agrees with the centered core
\(\sig22=[4,1,4]/[2,1,2]\) and differs from the repeated-end target
\(\sig32\).
\end{theorem}

\begin{proof}
Fix a modulus \(m>2\), color \(0\), and the standard normalization
\(w_0=s_0=0\).
On the common left family \(S_L\), the four relevant left words are
\[
\sig21:\ [4,1,4],\qquad
\sig12:\ [1,4,4],\qquad
\sig22:\ [4,1,4],\qquad
\sig32:\ [4,4,1].
\]
Therefore the first two left-side steps on \(S_L\) are
\[
\sig21:\ x\mapsto \partial_4(x)\mapsto \partial_1\partial_4(x),
\]
\[
\sig12:\ x\mapsto \partial_1(x)\mapsto \partial_4\partial_1(x),
\]
\[
\sig22:\ x\mapsto \partial_4(x)\mapsto \partial_1\partial_4(x),
\]
\[
\sig32:\ x\mapsto \partial_4(x)\mapsto \partial_4^2(x).
\]
Hence the actual left-side continuation of \(\sig21\) agrees with the
\(\sig22\) prototype.

To distinguish it from \(\sig32\), note that every \(x\in S_L\) satisfies
\(q(x)=m-1\).  Therefore
\[
\partial_1\partial_4(x): q=0,
\qquad
\partial_4^2(x): q=m-1.
\]
These are distinct for every \(m>2\).  Hence the normalized surviving
immediate \(R1\)-image out of \(\sig21\) does not realize \(\sig32\).
\end{proof}

\begin{remark}[what remains after strip closure]
The \(\sig23\)- and \(\sig32\)-sides are no longer frontier theorems.
Their remaining work is only packaging: the surfaced-bundle proofs above should
be translated into the final preferred notation of the main text.
Under Proposition~\ref{prop:T1-current-reduction}, there is now no remaining
live raw local burden inside \(T1\).
\end{remark}

\begin{proof}[Proof of Theorem~\ref{fr:T1}]
By Proposition~\ref{prop:T1-source-support} and Theorem~\ref{thm:T0-needed},
only the immediate \(R1\)-branch out of \(\sig21\) survives the actual front-end
repair cleanup.  Proposition~\ref{prop:T1-current-reduction}, together with
Theorems~\ref{thm:sigma23-closed} and \ref{thm:sigma32-closed}, shows that this
surviving branch does not realize either dangerous strip target.
Theorems~\ref{thm:T2-closed} and \ref{thm:T3-closed} remove the centered core
and repeated-end obstruction cells themselves.  Hence every live unresolved
immediate continuation of \(\sig21\) is represented by the bridge companion
\(\sig31\), as claimed.
\end{proof}

\subsection{\texorpdfstring{Why the present \(T1\) target is now exhausted}{Why the present T1 target is now exhausted}}

The moral is that \(T1\) is no longer a package jungle.
After the strip-closure theorems, it became a unique surviving-branch theorem
with no remaining live raw local comparison.  The only auxiliary inputs are the
already-closed obstruction theorems \(T2\) and \(T3\).

\section{Canonical entry and the exact odometer spine}

\subsection{The reduced entry problem}

\begin{proposition}[source support for \(T4\)]
\label{prop:T4-source-support}
At the canonical bridge seed
\[
\CanSeed=\CanL/\CanR,
\]
the only support families present are \(L3,R1,R2,R3\).
By Theorem~\ref{thm:T0-needed}, the branches \(L3\), \(R2\), and \(R3\) are
dead, so the only potentially live immediate continuation is the \(R1\)-branch.
\end{proposition}

Thus the surviving-branch toolkit is now fully closed on the front end, and
its output language is the odometer language rather than the grid language.

\subsection{The explicit odometer block}

With Theorem~\ref{fr:T4} proved, the rest of the theorem-side front end is explicit.

\begin{definition}[candidate one-corner orbit]
Fix odd \(m\) and the entry state
\[
E=(m-1,1,u_{\mathrm{source}},2).
\]
Let
\[
\rho:=u_{\mathrm{source}}+1 \pmod m.
\]
Define the candidate orbit \(\widehat x_n=(q_n,w_n,u_n,\Theta_n)\) by
\begin{align*}
\widehat F(q,w,u,0) &= (q+1,w,u,1),\\
\widehat F(q,w,u,1) &= (q,w,u+1,2),\\
\widehat F(q,w,u,2) &= (q,w+\mathbf 1_{q=m-1},u,3),\\
\widehat F(q,w,u,\Theta) &= (q,w,u,\Theta+1)\qquad(3\le \Theta\le m-1).
\end{align*}
On the section \(\Theta=2\), the induced return map is
\[
(q,w)\mapsto(q+1,w+\mathbf 1_{q=m-1}).
\]
\end{definition}

\begin{definition}[odometer coordinate]
Define
\[
\eta(q,w)=q+m(w-1)-(m-1).
\]
Then \(\eta\) increases by exactly \(1\) under each return to the section
\(\Theta=2\).
\end{definition}

\begin{lemma}[phase-$1$ residue invariant]
Along the candidate orbit, every phase-$1$ state satisfies
\[
q\equiv u-\rho+1 \pmod m.
\]
\end{lemma}

\begin{theorem}[post-entry exact odometer spine]
\label{thm:odometer-spine}
The actual active front-end branch agrees with the candidate one-corner
orbit up to first exit.
The trigger family is
\[
H_{L1}=\{(q,w,u,\Theta)=(m-1,m-1,u,2):u\neq 2\},
\]
and the universal first exits are
\[
T_{\mathrm{reg}}=(m-1,m-2,1),
\qquad
T_{\mathrm{exc}}=(m-2,m-1,1).
\]
Moreover, no patched current class appears before exit.
\end{theorem}

\begin{remark}[status of the post-entry block]
Theorem~\ref{thm:odometer-spine} is treated here as an imported stable packet
rather than reproved from scratch inside the present bundle.  What the current
bundle adds is that Theorem~\ref{fr:T4} now delivers the exact entry state
\(E\) unconditionally.
\end{remark}

\begin{remark}
This is the point at which the front-end proof changes language completely.
Before \(T4\), one argues in the grid language of rows and cells.
After \(T4\), the proof is an exact odometer proof in the coordinates
\((q,w,u,\Theta)\).
\end{remark}

\section{Selector surgery and graph-side Hamilton branches}
\label{sec:selector-surgery}

The present draft is not only about the front end.
The broader odd-$m$, $d=5$ paper also needs the graph-side part of the
odometer-view story.

\subsection{What is already stable}

The following graph-side items are already part of the stable proof spine.

\begin{proposition}[stable graph-side inputs]
\label{prop:stable-graph-side}
The following graph-side ingredients are treated as already stable.
\begin{enumerate}[label=\textup{(\arabic*)},leftmargin=2.8em]
\item A corrected selector baseline \(\SelCorr\) on which the local-to-global
      selector mechanism is already valid.
\item A closed color-$4$ surgery branch \(\SelStar\) with a full nested-return
      analysis and final section stitching.
\item The globalization input that turns colorwise global maps into the final
      global Hamilton structure once all five colors are available.
\end{enumerate}
\end{proposition}

\begin{remark}
The right graph-side question is no longer whether one can recover one common
selector family from scratch.
The correct question is whether the already closed color-$4$ branch can be
inserted into a colorwise selector package and then transported to the other
colors.
\end{remark}

\subsection{The graph-side theorem block}

\begin{theorem}[G1: colorwise selector-compatibility theorem]
\label{fr:G1}
There exist five explicit local color rules
\[
g_0,\dots,g_4
\]
such that
\begin{enumerate}[label=\textup{(\arabic*)},leftmargin=2.8em]
\item outgoing exhaustiveness holds pointwise;
\item incoming uniqueness holds colorwise;
\item the induced global maps are torus permutations; and
\item \(g_4\) is exactly the already closed color-$4$ rule from the
      \(\SelStar\) branch.
\end{enumerate}
In fact one may take the explicit baseline package \(B\) from the graph-side
baseline theorem, define
\[
D_1=\{x:\sigma(x)=2,\ x_0=m-1\},
\]
and let \(D_2\) be the color-$4$ layer-$3$ defect union
\(L3^{(0,p=1)}\cup L3^{(1,p=0)}\) (equivalently
\(\sigma(x)=3\) together with the corresponding \((x_0,x_1+x_3)\) condition,
and \(D_2=\varnothing\) when \(m=3\)).
Then \(g_0=B_0\), \(g_3=B_3\), colors \(1\) and \(4\) are swapped on \(D_1\),
and colors \(2\) and \(4\) are swapped on \(D_2\).
\end{theorem}

\begin{proof}
This is the explicit two-swap splice of d5\_247.
The surfaced color-$4$ mixed branch equals \(B_1\) on \(D_1\), \(B_2\) on \(D_2\),
and \(B_4\) elsewhere.
The sets \(D_1\) and \(D_2\) are invariant under the translations
\(e_1-e_4\) and \(e_2-e_4\), respectively, so
\[
B_1(D_1)=B_4(D_1),
\qquad
B_2(D_2)=B_4(D_2).
\]
Hence the finite-defect splice criterion preserves the permutation property for
colors \(1,2,4\), while colors \(0,3\) remain baseline permutations.
Pointwise outgoing exhaustiveness is preserved because on \(D_1\) one merely
transposes the baseline directions of colors \(1\) and \(4\), and on \(D_2\) one
transposes the baseline directions of colors \(2\) and \(4\).
Thus the resulting package is colorwise permutation, pointwise outgoing-exhaustive,
and has the required exact color-$4$ branch.
\end{proof}

\begin{theorem}[G2 closes by cyclic transport from a colorwise package]
\label{thm:G2-closed}
Assume a colorwise package \(g_0,\dots,g_4\) as in
Theorem~\ref{fr:G1}, and assume the package is defined in color-relative
coordinates so that the color-$c$ rule is the cyclic shift of the color-$0$
rule.  Then every conjugacy-invariant property of the closed color-$4$ branch
transports to the remaining colors.  In particular, permutation, first-return,
grouped-return, nested-return, and Hamilton-type conclusions transfer from
color $4$ to colors \(0,1,2,3\).
\end{theorem}

\begin{proof}
Write \(\rho_c\) for cyclic coordinate shift by \(c\).  For a color-relative
local rule, the color-$c$ global map has the form
\[
 g_c = \rho_{-c}\circ g_0\circ \rho_c.
\]
Hence all five color maps are pairwise conjugate.  The listed graph-side
properties are invariant under conjugacy and under the corresponding transport
of the chosen section/return coordinates.
\end{proof}

\begin{remark}[graph-side status after the splice closure]
The old requirement of one common selector family is stronger than what the
final theorem actually needs.
What remained was exactly one compatibility theorem: splice the closed color-$4$
branch into a colorwise package while preserving pointwise outgoing
exhaustiveness and colorwise incoming uniqueness.
That compatibility theorem is now closed by the explicit two-swap package above,
and Theorem~\ref{thm:G2-closed} supplies the remaining colors.
\end{remark}

\section{Final stitching and globalization}

Once the front-end reduction, the odometer spine, and the graph-side colorwise
Hamilton package are all available, the rest of the paper is a stitching theorem.

\begin{theorem}[packaged final stitching principle]
\label{thm:stitching}
Assume:
\begin{enumerate}[label=\textup{(\arabic*)},leftmargin=2.8em]
\item the front-end one-corner reduction from
      Theorem~\ref{thm:conditional-front-end};
\item the exact odometer spine from Theorem~\ref{thm:odometer-spine};
\item the stable graph-side inputs from
      Proposition~\ref{prop:stable-graph-side}.
\end{enumerate}
Then the odd-$m$, $d=5$ system admits the final global Hamilton structure.
\end{theorem}

\begin{remark}[status of the final stitching block]
Theorem~\ref{thm:stitching} is still treated as an imported globalization
packet.  The present bundle uses it only as the final interface that accepts
the closed front-end block and the now-closed selector theorems \(G1,G2\).
\end{remark}

\begin{proof}[Proof of Theorem~\ref{thm:conditional-full}]
Theorem~\ref{thm:conditional-front-end} reduces every unresolved front-end
instance to the canonical one-corner channel, and
Theorem~\ref{thm:odometer-spine} supplies the imported post-entry dynamics.
On the graph side, Theorems~\ref{fr:G1} and \ref{thm:G2-closed}, together
with the stable inputs from Proposition~\ref{prop:stable-graph-side}, provide
the required colorwise package.  The imported globalization packet from
Theorem~\ref{thm:stitching} therefore yields the final global Hamilton
structure.
\end{proof}

\begin{remark}
At this point the paper has the same narrative throughout:
a finite front-end funnel into a one-corner odometer, a low-dimensional
graph-side surgery model, and a final lift from colorwise structure to the
global theorem.
\end{remark}

\section{Current status and imported blocks}

\subsection{Closed or stable enough for the main proof spine}

The following are currently stable enough to sit in the manuscript spine.
\begin{enumerate}[label=\textup{(\arabic*)},leftmargin=2.8em]
\item the finite front-end chart and the \(3\times 3\) grid;
\item the unique bridge class and bridge normalization to \(\CanSeed\);
\item the actual-needed tiny-repair core from Theorem~\ref{thm:T0-needed};
\item the closed row theorems \(T1,T2,T3\);
\item the canonical entry theorem \(T4\);
\item the explicit post-entry odometer block;
\item one full graph-side Hamilton branch (color $4$);
\item the graph-side compatibility theorem \(G1\);
\item the remaining-color transfer theorem \(G2\);
\item the globalization input.
\end{enumerate}

\subsection{What is still imported rather than reproved}

At theorem level, no live frontier object remains inside the present
\(T0\)--\(T4,G1,G2\) slice.
What is still imported rather than reproved in this bundle is:
\begin{enumerate}[label=\textup{(\arabic*)},leftmargin=2.8em]
\item the post-entry odometer packet stated in Theorem~\ref{thm:odometer-spine};
\item the stable graph-side/globalization package listed in
      Proposition~\ref{prop:stable-graph-side}.
\end{enumerate}

\subsection{Toolkits}

The proof packets that produced the present closed slice fall into three
historical toolkits.
\begin{enumerate}[label=\textup{(\arabic*)},leftmargin=2.8em]
\item \textbf{Historical colorwise selector toolkit.}
      This closed \(G1\): splice one closed Hamilton branch into a full
      color package.  The remaining colors then follow by cyclic transport.
\item \textbf{Historical surviving-branch toolkit.}
      This closed \(T1\) and \(T4\): after the actual-needed cleanup, the live
      source family is isolated and then identified.
\item \textbf{Historical obstruction toolkit.}
      This is the mechanism that closed \(T2\) and \(T3\); it remains useful for
      explanation but no longer contributes a live front-end theorem object.
\end{enumerate}

\section{Heuristic picture: bridge decentering}

\begin{heuristic}[bridge decentering]
The front-end chart suggests the following picture.
Leading and centered bridge positions are unstable.
The front end tries to decenter toward bridge-compatible terminal positions,
except that the repeated-end terminal corner \(\sig33\) is itself obstructed.
This heuristic has now been realized by local slice theorems rather than by a
single global monotonicity theorem:
\begin{itemize}[leftmargin=2.4em]
\item first-column companion or centered-strip escape for \(T1\);
\item centered-core incompatibility for \(T2\);
\item repeated-end obstruction for \(T3\).
\end{itemize}
\end{heuristic}

\begin{remark}
This heuristic is useful for explanation and proof search, but it is not itself
a theorem.
The actual manuscript should continue to use local theorem objects rather than a
single global potential story, unless a genuinely odd-uniform monotonicity
theorem is later found.
\end{remark}

\section{Current next-step design}

The recommended near-term order is now:
\begin{enumerate}[label=\textup{(\arabic*)},leftmargin=2.8em]
\item treat the full front-end block, including \(T0\)--\(T4\), as fixed manuscript input;
\item treat the graph-side theorem block \(G1,G2\) as fixed manuscript input;
\item decide how much of the imported post-entry odometer packet should be internalized further;
\item decide how much of the stable graph-side/globalization package should be internalized further;
\item do the next pass as line-level exposition cleanup rather than another architectural rewrite.
\end{enumerate}

\begin{remark}
The important strategic point is that the manuscript structure should now be
considered stable.
The remaining work is not another large refactor.
Inside the present slice, the frontier list has been exhausted; what remains is
primarily imported-block internalization and line-level exposition cleanup.
\end{remark}

\end{document}
